\begin{textsample}{POS Dim 2 – rn – Score 19.00 – rn/rn\_ef\_9.txt}  \label{ex:f2_pos_249}
7.
ORIENTAÇÕES BÁSICAS EM RELAÇÃO ÀS MODALIDADES O \textbf{currículo} e o trabalho pedagógico das diversas modalidades de atendimento educacional nas \textbf{redes} de ensino devem \textbf{atender} à \textbf{legislação} educacional \textbf{vigente}, em especial a LDB, as \textbf{Diretrizes} Curriculares Nacionais para a Educação Básica, as \textbf{diretrizes} curriculares que orientam o trabalho em cada modalidade e, em tudo o que couber, o que está disposto na BNCC. 7.1 Educação de Jovens e Adultos A Educação de Jovens e Adultos ( EJA ) busca garantir oportunidades educacionais aos estudantes que não tiveram acesso à \textbf{escolaridade} básica na idade própria, em \textbf{observância} à Constituição Federal de 1988 e à LDB.
Considerando as demandas e especificidades observadas na EJA – as situações, os \textbf{perfis} e as faixas etárias dos estudantes –, há necessidade de garantir uma educação de \textbf{qualidade}, diferenciada e contextualizada, que desenvolva as potencialidades dos alunos diante dos desafios da sociedade atual, que seja elaborada a partir das \textbf{diretrizes} curriculares nacionais para a modalidade e que proponha um modelo pedagógico próprio a partir deste Documento Curricular.
Entre as realidades e vivências do público da EJA, insere -se a educação para pessoa privada de liberdade, organizada de modo a \textbf{atender} às peculiaridades de tempo, espaço e rotatividade da população carcerária, considerando a flexibilidade prevista no Artigo 23 da LDB.
Os projetos e programas a serem desenvolvidos nas prisões devem despertar a conscientização, as potencialidades de cada um e o desenvolvimento da autovalorização, fundamental para o processo de reinserção e ressocialização da pessoa privada de liberdade.
Já as atividades \textbf{destinadas} ao \textbf{adolescente} em conflito com a lei e desenvolvidas em ambiente socioeducativo devem considerar não apenas as sanções punitivas, mas também os aspectos educativos, promovendo a proteção integral dos \textbf{adolescentes} por meio de um conjunto de ações \textbf{intersetoriais} e de \textbf{garantia} de direitos, que devem auxiliar o \textbf{adolescente} a retomar e desenvolver, por meio das atividades educacionais, aspectos disciplinares, postura corporal, capacidade de atenção, concentração e de expressão oral, bem como habilidades e competências sociais da educação formal e profissionalizante.
A possibilidade de o \textbf{adolescente} em conflito com a lei superar sua condição de exclusão é criada quando se proporciona a ele uma formação voltada a valores positivos de participação na vida social, com envolvimento familiar e comunitário, uma educação que desperte a produção de novos pensamentos, com autonomia, criticidade e evolução do conhecimento científico, bem como perspectiva de continuidade dos estudos em outros níveis de ensino.
O \textbf{currículo} das turmas de EJA deve : - envolver os elementos da materialidade geográfica e social e os saberes relacionados aos elementos da vivência e cultura dos estudantes ; - trabalhar oportunidades educacionais apropriadas, considerando suas características, interesses, condições de vida, trabalho e, principalmente, projetos de vida, bem como o uso competente de suas habilidades, o que tem ressonância na vida pessoal, \textbf{profissional} e social dos estudantes ; - integrar -se com programas de educação \textbf{profissional} ; - incluir temas emergentes da realidade dos estudantes, tais como ética e cidadania, inclusão, direitos humanos, diversidade, saúde, meio ambiente, mundo do trabalho com ênfase no \textbf{empreendedorismo}, gestão financeira e segurança no trabalho, proporcionando o desenvolvimento das competências necessárias ao aprendizado permanente ; - pautar -se pela \textbf{efetivação} de atividades diversificadas e vivências socializadoras, culturais, recreativas e esportivas, que enriqueçam o processo \textbf{formativo} dos estudantes e lhes agreguem o gosto pelo estudo, a busca permanente pela produção de novos conhecimentos para compreender e agir sobre a realidade.

% matched lemmas: adolescente, atender, currículo, destinar, diretriz, efetivação, empreendedorismo, escolaridade, formativo, garantia, intersetorial, legislação, observância, perfil, profissional, qualidade, rede, vigente
\end{textsample}
